% ---------------------------------------------------------
% TABELA EMENTA 37: Gestão de Marketing I
% ---------------------------------------------------------
\begin{xltabular}{\linewidth}{|X|p{2.5cm}|p{2.5cm}|}
\hline
\multirow{3}{=}{\textcolor{ifscgreen}{\textbf{Unidade Curricular:}}\\[3pt]\textbf{\large Gestão de Marketing I}} & \multicolumn{2}{p{\dimexpr5cm+2\tabcolsep+\arrayrulewidth\relax}|}{\textcolor{ifscgreen}{\textbf{Semestre:}} \textbf{2º (Ano 1)}} \\ \cline{2-3}
 & \textcolor{ifscgreen}{\textbf{CH EaD*:}} & \textcolor{ifscgreen}{\textbf{CH Total*:}} \\ \cline{2-3}
 & \textbf{00 h} & \textbf{40 h} \\ \hline
\endhead
\multicolumn{3}{|p{\dimexpr\linewidth-2\tabcolsep-2\arrayrulewidth\relax}|}{\textcolor{ifscgreen}{\textbf{Objetivos:}}} \\ \hline
\multicolumn{3}{|p{\dimexpr\linewidth-2\tabcolsep-2\arrayrulewidth\relax}|}{
\begin{itemize}[leftmargin=*,noitemsep,topsep=2pt]
 \item \textit{[Objetivos de aprendizagem pendentes de preenchimento pelo docente responsável]}
\end{itemize}
} \\ \hline
\multicolumn{3}{|p{\dimexpr\linewidth-2\tabcolsep-2\arrayrulewidth\relax}|}{\textcolor{ifscgreen}{\textbf{Conteúdos:}}} \\ \hline
\multicolumn{3}{|p{\dimexpr\linewidth-2\tabcolsep-2\arrayrulewidth\relax}|}{
\begin{itemize}[leftmargin=*,noitemsep,topsep=2pt]
 \item \textit{[Conteúdo programático pendente de preenchimento pelo docente responsável]}
\end{itemize}
} \\ \hline
\multicolumn{3}{|p{\dimexpr\linewidth-2\tabcolsep-2\arrayrulewidth\relax}|}{\textcolor{ifscgreen}{\textbf{Estratégias de Ensino e Aprendizagem:}}} \\ \hline
\multicolumn{3}{|p{\dimexpr\linewidth-2\tabcolsep-2\arrayrulewidth\relax}|}{\textit{[Metodologia de ensino e avaliação pendente de preenchimento pelo docente responsável]}} \\ \hline
\multicolumn{3}{|p{\dimexpr\linewidth-2\tabcolsep-2\arrayrulewidth\relax}|}{\textcolor{ifscgreen}{\textbf{Bibliografia Básica:}}} \\ \hline
\multicolumn{3}{|p{\dimexpr\linewidth-2\tabcolsep-2\arrayrulewidth\relax}|}{1. \textit{[1º título da Bibliografia Básica pendente de indicação docente e validação no Parecer da Biblioteca do Câmpus]}
2. \textit{[2º título da Bibliografia Básica pendente de indicação docente e validação no Parecer da Biblioteca do Câmpus]}} \\ \hline
\multicolumn{3}{|p{\dimexpr\linewidth-2\tabcolsep-2\arrayrulewidth\relax}|}{\textcolor{ifscgreen}{\textbf{Bibliografia Complementar:}}} \\ \hline
\multicolumn{3}{|p{\dimexpr\linewidth-2\tabcolsep-2\arrayrulewidth\relax}|}{1. \textit{[1º título da Bibliografia Complementar pendente de indicação docente e validação no Parecer da Biblioteca do Câmpus]}
2. \textit{[2º título da Bibliografia Complementar pendente de indicação docente e validação no Parecer da Biblioteca do Câmpus]}
3. \textit{[3º título da Bibliografia Complementar pendente de indicação docente e validação no Parecer da Biblioteca do Câmpus]}} \\ \hline
\end{xltabular}

\vspace{0.4cm}


% ---------------------------------------------------------
% TABELA EMENTA 38: Organização e Processos
% ---------------------------------------------------------
\begin{xltabular}{\linewidth}{|X|p{2.5cm}|p{2.5cm}|}
\hline
\multirow{3}{=}{\textcolor{ifscgreen}{\textbf{Unidade Curricular:}}\\[3pt]\textbf{\large Organização e Processos}} & \multicolumn{2}{p{\dimexpr5cm+2\tabcolsep+\arrayrulewidth\relax}|}{\textcolor{ifscgreen}{\textbf{Semestre:}} \textbf{2º (Ano 1)}} \\ \cline{2-3}
 & \textcolor{ifscgreen}{\textbf{CH EaD*:}} & \textcolor{ifscgreen}{\textbf{CH Total*:}} \\ \cline{2-3}
 & \textbf{00 h} & \textbf{40 h} \\ \hline
\endhead
\multicolumn{3}{|p{\dimexpr\linewidth-2\tabcolsep-2\arrayrulewidth\relax}|}{\textcolor{ifscgreen}{\textbf{Objetivos:}}} \\ \hline
\multicolumn{3}{|p{\dimexpr\linewidth-2\tabcolsep-2\arrayrulewidth\relax}|}{
\begin{itemize}[leftmargin=*,noitemsep,topsep=2pt]
 \item Analisar as estruturas organizacionais e sua influência na organização do trabalho e dos processos administrativos.
 \item Identificar os elementos que compõem um processo organizacional, compreendendo a relação entre atividades, pessoas, recursos e informações.
 \item Representar processos administrativos utilizando ferramentas apropriadas, como fluxogramas e outros instrumentos de mapeamento.
 \item Auxiliar na análise de processos organizacionais, identificando oportunidades de melhoria relacionadas à eficiência, à comunicação e à padronização das atividades.
 \item Aplicar princípios de organização administrativa na elaboração e padronização de documentos, procedimentos e rotinas de trabalho.
\end{itemize}
} \\ \hline
\multicolumn{3}{|p{\dimexpr\linewidth-2\tabcolsep-2\arrayrulewidth\relax}|}{\textcolor{ifscgreen}{\textbf{Conteúdos:}}} \\ \hline
\multicolumn{3}{|p{\dimexpr\linewidth-2\tabcolsep-2\arrayrulewidth\relax}|}{
\begin{itemize}[leftmargin=*,noitemsep,topsep=2pt]
 \item Organização do trabalho nas organizações: conceitos e características.
 \item Estruturas organizacionais: tipos, departamentalização e formas de delegação, centralização e descentralização.
 \item Organização do ambiente de trabalho: layout, fluxo e comunicação entre pessoas e setores.
 \item Processos administrativos: conceito, identificação, mapeamento, ferramentas de representação e análise.
 \item Organização, padronização e documentação de procedimentos, formulários e manuais administrativos.
\end{itemize}
} \\ \hline
\multicolumn{3}{|p{\dimexpr\linewidth-2\tabcolsep-2\arrayrulewidth\relax}|}{\textcolor{ifscgreen}{\textbf{Estratégias de Ensino e Aprendizagem:}}} \\ \hline
\multicolumn{3}{|p{\dimexpr\linewidth-2\tabcolsep-2\arrayrulewidth\relax}|}{- **Metodologia:** A unidade curricular será desenvolvida por meio de aulas expositivas dialogadas e metodologias ativas de aprendizagem, com enfoque teórico-prático, favorecendo a compreensão da organização do trabalho e dos processos administrativos nas organizações. As estratégias de ensino buscarão aproximar os estudantes de situações reais do contexto organizacional, promovendo a observação, representação, análise e proposição de melhorias em processos e rotinas administrativas. Entre as estratégias metodológicas poderão ser utilizados estudos de caso, resolução de situações-problema, mapeamento e representação de processos por meio de fluxogramas e outras ferramentas, análise de documentos e procedimentos organizacionais, simulações, atividades práticas, pesquisas, trabalhos colaborativos, visitas técnicas e interações com organizações e profissionais da área. Por conter ferramentas e métodos relacionados à organização, diagnóstico, representação e análise, esta unidade curricular se articula com Introdução à Administração, Oficina de Integração I e II, Gestão de Pessoas e Gestão de Operações e Qualidade, além de dialogar com a formação geral (Língua Portuguesa, Matemática, Sociologia, Sociedade e Trabalho, História e Geografia).
- **Avaliação:** Contínua, processual e formativa, contemplando diferentes instrumentos, como produções escritas, apresentações orais, resolução de estudos de caso, atividades práticas, relatórios, registros reflexivos e participação nas atividades propostas, individualmente ou em grupo.} \\ \hline
\multicolumn{3}{|p{\dimexpr\linewidth-2\tabcolsep-2\arrayrulewidth\relax}|}{\textcolor{ifscgreen}{\textbf{Bibliografia Básica:}}} \\ \hline
\multicolumn{3}{|p{\dimexpr\linewidth-2\tabcolsep-2\arrayrulewidth\relax}|}{1. ARAÚJO, Luiz Cesar G. de. \textbf{Organização, sistemas e métodos}. São Paulo: Atlas, 2010.
2. CRUZ, Tadeu. \textbf{Processos organizacionais e métodos}. 5. ed. São Paulo: Atlas, 2021.
3. BERNARDI, Luiz Antonio. \textbf{Manual de empreendedorismo e gestão: fundamentos, estratégias e dinâmicas}. 2. ed. São Paulo: Atlas, 2012.} \\ \hline
\multicolumn{3}{|p{\dimexpr\linewidth-2\tabcolsep-2\arrayrulewidth\relax}|}{\textcolor{ifscgreen}{\textbf{Bibliografia Complementar:}}} \\ \hline
\multicolumn{3}{|p{\dimexpr\linewidth-2\tabcolsep-2\arrayrulewidth\relax}|}{1. CHIAVENATO, Idalberto. \textbf{Administração para não administradores: a gestão de negócios ao alcance de todos}. 2. ed. rev. e atual. Barueri: Manole, 2011.
2. CINTRA, M.; CUNHA, M. P. \textbf{Rotinas administrativas}. São Paulo: KCM, 2008.
3. FARIAS, Cláudio V. S. (org.). \textbf{Técnico em administração: gestão e negócios}. Porto Alegre: Bookman, 2013.} \\ \hline
\end{xltabular}

\vspace{0.4cm}


% ---------------------------------------------------------
% TABELA EMENTA 39: Informática Aplicada
% ---------------------------------------------------------
\begin{xltabular}{\linewidth}{|X|p{2.5cm}|p{2.5cm}|}
\hline
\multirow{3}{=}{\textcolor{ifscgreen}{\textbf{Unidade Curricular:}}\\[3pt]\textbf{\large Informática Aplicada}} & \multicolumn{2}{p{\dimexpr5cm+2\tabcolsep+\arrayrulewidth\relax}|}{\textcolor{ifscgreen}{\textbf{Semestre:}} \textbf{1º (Ano 1)}} \\ \cline{2-3}
 & \textcolor{ifscgreen}{\textbf{CH EaD*:}} & \textcolor{ifscgreen}{\textbf{CH Total*:}} \\ \cline{2-3}
 & \textbf{00 h} & \textbf{40 h} \\ \hline
\endhead
\multicolumn{3}{|p{\dimexpr\linewidth-2\tabcolsep-2\arrayrulewidth\relax}|}{\textcolor{ifscgreen}{\textbf{Objetivos:}}} \\ \hline
\multicolumn{3}{|p{\dimexpr\linewidth-2\tabcolsep-2\arrayrulewidth\relax}|}{
\begin{itemize}[leftmargin=*,noitemsep,topsep=2pt]
 \item Reconhecer componentes do computador e sistemas operacionais.
 \item Compreender segurança da informação e cidadania digital.
 \item Organizar arquivos e estruturar informações digitais.
 \item Aplicar recursos de processadores de texto para elaborar documentos administrativos profissionais (cartas, relatórios, memorandos, ofícios).
 \item Aplicar fórmulas e funções em planilhas eletrônicas para organização e análise de dados.
 \item Criar gráficos e tabelas dinâmicas para visualização de informações.
 \item Aplicar técnicas de design visual em apresentações de slides profissionais.
 \item Aplicar ferramentas de computação em nuvem para colaboração e compartilhamento de documentos.
 \item Aplicar ferramentas de IA generativa de forma responsável e ética no apoio a estudos e atividades administrativas.
\end{itemize}
} \\ \hline
\multicolumn{3}{|p{\dimexpr\linewidth-2\tabcolsep-2\arrayrulewidth\relax}|}{\textcolor{ifscgreen}{\textbf{Conteúdos:}}} \\ \hline
\multicolumn{3}{|p{\dimexpr\linewidth-2\tabcolsep-2\arrayrulewidth\relax}|}{
\begin{itemize}[leftmargin=*,noitemsep,topsep=2pt]
 \item Fundamentos de Informática: hardware, software e sistemas operacionais; componentes básicos do computador; computação em nuvem e colaboração digital.
 \item Segurança da informação e cidadania digital.
 \item Organização de arquivos e pastas; edição de imagens.
 \item Editores de texto: formatação de documentos, criação de tabelas e inserção de imagens.
 \item Planilhas eletrônicas: funcionalidade e navegação, fórmulas e funções, formatação condicional, gráficos, análise de dados e tabelas dinâmicas.
 \item Ferramentas de apresentação profissional: criação e personalização de apresentações; formulários.
 \item Inteligência Artificial Generativa: conceitos básicos, ferramentas de IA, suporte a pesquisa e aprendizagem, uso ético e responsável.
 \item Temas transversais: segurança em redes sociais, privacidade de dados e combate à desinformação/fake news (educação em direitos humanos e prevenção de violência virtual); ergonomia e saúde na postura durante o uso prolongado de telas; descarte responsável de equipamentos eletrônicos e gestão de resíduos (educação ambiental); atuação crítica e responsável com ferramentas de IA generativa.
\end{itemize}
} \\ \hline
\multicolumn{3}{|p{\dimexpr\linewidth-2\tabcolsep-2\arrayrulewidth\relax}|}{\textcolor{ifscgreen}{\textbf{Estratégias de Ensino e Aprendizagem:}}} \\ \hline
\multicolumn{3}{|p{\dimexpr\linewidth-2\tabcolsep-2\arrayrulewidth\relax}|}{- **Metodologia:** Desenvolvimento predominantemente prático, com atividades no Laboratório de Informática (carga horária prática prevista de 30h e teórica de 10h), utilizando slides e quadro branco para a exposição dialogada. Articula-se com Comunicação e Linguagem, Projeto Integrador e Administração Financeira/Marketing. Recursos necessários: laboratório de informática com computadores conectados à internet, softwares de automação de escritório (LibreOffice ou Microsoft Office), navegadores web, ferramentas de colaboração em nuvem (Google Workspace, Microsoft 365), projetor, quadro branco e acesso responsável a ferramentas de IA generativa. Critério de divisão de turma: turmas com mais de 28 alunos são divididas entre dois docentes atuando em paralelo em dois laboratórios diferentes, devido à lotação máxima dos laboratórios de informática.
- **Avaliação:** Contempla trabalhos em grupo, resolução de exercícios práticos e produção de atividades integradoras que relacionem teoria e prática em situações reais de uso das ferramentas digitais.} \\ \hline
\multicolumn{3}{|p{\dimexpr\linewidth-2\tabcolsep-2\arrayrulewidth\relax}|}{\textcolor{ifscgreen}{\textbf{Bibliografia Básica:}}} \\ \hline
\multicolumn{3}{|p{\dimexpr\linewidth-2\tabcolsep-2\arrayrulewidth\relax}|}{1. CAPRON, H. L.; JOHNSON, J. A. \textbf{Introdução à informática}. Tradução de José Carlos Barbosa dos Santos. 8. ed. São Paulo: Pearson Prentice Hall, 2004.
2. MANZANO, André Luiz N. G.; MANZANO, Maria Izabel N. G. \textbf{Estudo dirigido de informática básica}. 7. ed. atual., rev. e ampl. São Paulo: Érica, 2007.} \\ \hline
\multicolumn{3}{|p{\dimexpr\linewidth-2\tabcolsep-2\arrayrulewidth\relax}|}{\textcolor{ifscgreen}{\textbf{Bibliografia Complementar:}}} \\ \hline
\multicolumn{3}{|p{\dimexpr\linewidth-2\tabcolsep-2\arrayrulewidth\relax}|}{1. CERT.BR. \textbf{Cartilha de Segurança para Internet: fascículos}. [S. l.]: CERT.br, [s. d.]. Disponível em: https://cartilha.cert.br/fasciculos/.
2. JR., Edgard Bruno C. \textbf{Informática Aplicada às Áreas de Contabilidade, Administração e Economia}. 4. ed. Rio de Janeiro: Atlas, 2012. E-book. ISBN 9788522494651.
3. LIBREOFFICE. \textbf{Guia do LibreOffice}. [s.l.]: The Document Foundation, 2025. Disponível em: https://books.libreoffice.org/pt-br/.
4. MARTY, Nicolas. \textbf{Inteligência artificial - O futuro é agora: dialogar com a IA}. [S.l.]: Publicação independente, 2023. ISBN 9798857441145.
5. OLIVEIRA, Rômulo Silva de; CARISSIMI, Alexandre da Silva; TOSCANI, Simão Sirineo. \textbf{Sistemas operacionais}. 4. ed. Porto Alegre: Bookman: Instituto de Informática da UFRGS, 2010.
6. SILVA, Mário Gomes da. \textbf{Informática - Terminologia - Microsoft Windows 8 - Internet - Segurança - Microsoft Word 2013 - Microsoft Excel 2013 - Microsoft PowerPoint 2013 - Microsoft Access 2013}. Rio de Janeiro: Érica, 2013. E-book. ISBN 9788536519319.} \\ \hline
\end{xltabular}

\vspace{0.4cm}

